\subsubsection{从米勒指数找到法向}

晶面法向由倒格矢确定。采用第\ref{ch:theory}章的倒格约定，表面对应的方向为
\begin{equation}
  \mathbf G_{hkl}=h\mathbf b_1+k\mathbf b_2+l\mathbf b_3,
  \qquad
  \hat{\mathbf n}=\frac{\mathbf G_{hkl}}{\|\mathbf G_{hkl}\|}.
  \label{eq:surface-miller-normal}
\end{equation}
式中，$\mathbf b_1,\mathbf b_2,\mathbf b_3$为倒格矢，单位为$\text{\AA}^{-1}$；$h,k,l$是无量纲整数，$\hat{\mathbf n}$是单位法向量。立方晶胞的三条边等长且正交，$(111)$法向与$\mathbf a+\mathbf b+\mathbf c$平行；本例的晶格轴沿全局坐标轴排列，因此也与全局$(1,1,1)$平行。斜晶胞一般没有这种对应关系，不能把米勒指数直接当作笛卡尔方向输入吸附物的定向参数。

构建后的表面晶胞以$\mathbf a,\mathbf b$为面内矢量，可用
\begin{equation}
  \hat{\mathbf n}_{\mathrm{slab}}=
  \frac{\mathbf a\times\mathbf b}{\|\mathbf a\times\mathbf b\|},
  \qquad H=|\mathbf c\cdot\hat{\mathbf n}_{\mathrm{slab}}|
  \label{eq:surface-normal-height}
\end{equation}
检查实际法向与晶胞高度。$H$是相邻周期表面之间沿法向的重复长度，并不总等于$\|\mathbf c\|$。把同一薄片的原子放在连续的周期映像中后，将原子位置投影到该法向，便可读出最外两层之间的跨度。它与构建窗口的最小厚度、完整重复单元厚度也应分别记录。

\begin{table}[H]
\centering
\caption{表面构建中几种长度的含义。}
\begin{tabularx}{\linewidth}{lX}
\toprule
量 & 怎样理解与核对\\
\midrule
最小薄片厚度 & 构建时用于选择完整重复单元的条件；不是直接指定层数。\\
原子层跨度 & 最外两层的法向距离；检查时先选取同一连续薄片。\\
法向胞高$H$ & 由式\eqref{eq:surface-normal-height}得到，包含薄片和周期间隔。\\
新增真空$\Delta$ & \ui{添加真空层}中沿第三晶格矢量增加的长度，不是最终总真空。\\
\bottomrule
\end{tabularx}
\end{table}

对于非正交输入，还要检查实际输出的法向、层数和周期间隔。当前工作构建对斜晶胞的厚度换算与候选约简可能受到取向处理的影响，不宜只依据输入的最小厚度判断结果。本章的Pt例子从已核对的立方晶胞开始，再检查输出表面。
